\section{Original ambient projective sections are original homogeneous polynomials}
\label{sec:c0n-ambient-section-polynomials}

Let $R$ be an original commutative ring, $P=\mathbb P_R^N$ its original
polynomial Proj, and let $e\geq0$ be integral. All coefficients below
belong to $R$ itself. No field of fractions or field-only global-section
representation is substituted for the original affine base. The target
line is the original coordinate-hyperplane twist, with its native frames.
A universe lift is applied only to the coordinate index of the auxiliary
Laurent algebra; literal original polynomial renaming and its inverse
recover the original coordinate type and original ring coefficients.

\begin{lemma}[The actual native coordinate fractions have their genuine Laurent comparison]
\label{lem:c0n-native-coordinate-laurent-comparison}
\lean{DeJongAlterations.Geometry.ProperSteinProjectiveCech.polynomialCoordinateHomogeneousLaurentEquiv}
\leanok
Let $R$ be an original commutative ring, $I$ a finite linearly ordered
coordinate set, $m\in\mathbf Z$ and $S\subset I$.  The actual homogeneous
degree-$m$ piece of localization at $\prod_{i\in S}x_i$ is linearly
equivalent to the entire actual Laurent module whose supported integral
exponents sum to $m$ and can be negative only at coordinates in $S$.
The map retains each original numerator and its original denominator.
\end{lemma}
\begin{proof}
Use the already-proved native comparison in
\nolinkurl{DeJongAlterations.Geometry.ProperSteinProjectiveCechPolynomialLocalization}.
Distinct original polynomial exponent vectors remain distinct integral
Laurent vectors, so the original polynomial map is injective.  Each
inverted coordinate is a unit, and the actual localization equality
criterion therefore proves injectivity of the genuine fraction map.
For surjectivity, take one supported Laurent monomial.  Choose a natural
power of the same coordinate product clearing its negative exponents.
The resulting exponents are all nonnegative and have precisely the
homogeneous degree of the required numerator.  The original polynomial
monomial divided by that product power has exactly the original Laurent
image.  Finite sums give every supported Laurent element.  The actual
linear map is consequently bijective, giving the native linear
equivalence without a field replacement or comparison certificate.
\end{proof}

\begin{lemma}[A whole nonnegative Laurent piece is a genuine polynomial piece]
\label{lem:c0n-ambient-empty-laurent-polynomial}
\leanok
\lean{EnvelopeProperty.Char0.RelativeProjectiveSectionPolynomial.integral_homogeneous_isHomogeneous,EnvelopeProperty.Char0.RelativeProjectiveSectionPolynomial.exists_polynomial_of_empty_laurent,EnvelopeProperty.Char0.RelativeProjectiveSectionPolynomial.polynomial_unique_of_same_laurent}
For a finite original coordinate type, every Laurent polynomial whose
supported exponents are nonnegative and have total degree $e$ is the
image of a unique original homogeneous polynomial of degree $e$.
\end{lemma}
\begin{proof}
Use the actual polynomial-to-Laurent inclusion, which is injective
because distinct natural exponent vectors remain distinct integral
exponent vectors. The native homogeneous coordinate-product localization
comparison for the empty coordinate set is onto this very Laurent
piece. Its inverse supplies a genuine fraction with original polynomial
numerator, total numerator degree $e+n\cdot0=e$, and denominator
$(\prod_{i\in\varnothing}x_i)^n=1$. The actual denominator-clearing
identity therefore identifies that same numerator's Laurent image with
the original element. Its integral weighted degree equals the natural
total degree, since every original weight is one. Injectivity gives
uniqueness. This proof works for arbitrary coefficient rings, including
the zero ring; it does not use a zero homogeneous prime.
\uses{lem:c0n-native-coordinate-laurent-comparison}
\end{proof}

\begin{lemma}[The same native global section has one compatible Laurent expression]
\label{lem:c0n-ambient-native-section-laurent}
\leanok
\lean{EnvelopeProperty.Char0.RelativeProjectiveSectionPolynomial.coefficient_transition_of_open_transition,EnvelopeProperty.Char0.RelativeProjectiveSectionPolynomial.nativeCoefficient,EnvelopeProperty.Char0.RelativeProjectiveSectionPolynomial.nativeCoefficient_overlap,EnvelopeProperty.Char0.RelativeProjectiveSectionPolynomial.nativeCoefficient_ext,EnvelopeProperty.Char0.RelativeProjectiveSectionPolynomial.nativeCoefficient_homogeneousSection,EnvelopeProperty.Char0.RelativeProjectiveSectionPolynomial.originalLaurent,EnvelopeProperty.Char0.RelativeProjectiveSectionPolynomial.originalPolynomialToLaurent,EnvelopeProperty.Char0.RelativeProjectiveSectionPolynomial.originalPolynomialToLaurent_X_isUnit,EnvelopeProperty.Char0.RelativeProjectiveSectionPolynomial.originalLocalizationToLaurent,EnvelopeProperty.Char0.RelativeProjectiveSectionPolynomial.originalLocalizationToLaurent_algebraMap,EnvelopeProperty.Char0.RelativeProjectiveSectionPolynomial.originalAwayToLaurent,EnvelopeProperty.Char0.RelativeProjectiveSectionPolynomial.originalAwayToLaurent_mk_clear,EnvelopeProperty.Char0.RelativeProjectiveSectionPolynomial.originalAwayToLaurent_awayMap,EnvelopeProperty.Char0.RelativeProjectiveSectionPolynomial.chartAwayEquiv,EnvelopeProperty.Char0.RelativeProjectiveSectionPolynomial.overlapAwayEquiv,EnvelopeProperty.Char0.RelativeProjectiveSectionPolynomial.chartToLaurent,EnvelopeProperty.Char0.RelativeProjectiveSectionPolynomial.overlapToLaurent,EnvelopeProperty.Char0.RelativeProjectiveSectionPolynomial.chartAway_restrict_left,EnvelopeProperty.Char0.RelativeProjectiveSectionPolynomial.chartAway_restrict_right,EnvelopeProperty.Char0.RelativeProjectiveSectionPolynomial.overlapToLaurent_restrict_left,EnvelopeProperty.Char0.RelativeProjectiveSectionPolynomial.overlapToLaurent_restrict_right,EnvelopeProperty.Char0.RelativeProjectiveSectionPolynomial.chartToLaurent_coefficient_clear,EnvelopeProperty.Char0.RelativeProjectiveSectionPolynomial.overlapToLaurent_transition_clear,EnvelopeProperty.Char0.RelativeProjectiveSectionPolynomial.originalAwayToLaurent_shift_mem,EnvelopeProperty.Char0.RelativeProjectiveSectionPolynomial.sectionLaurent,EnvelopeProperty.Char0.RelativeProjectiveSectionPolynomial.sectionLaurent_eq,EnvelopeProperty.Char0.RelativeProjectiveSectionPolynomial.sectionLaurent_mem_singleton}
For an original global section $s$ of the actual twist $\mathcal O(e)$,
let $c_i$ be its coefficient in the original frame on $D_+(x_i)$.
Map the genuine degree-zero homogeneous fraction $c_i$ into the original
Laurent algebra and multiply by $x_i^e$. The resulting expressions
$\lambda_i$ are equal for all original indices. Each has total degree
$e$, and only coordinate $i$ can have a negative exponent in
$\lambda_i$. The same native frame coefficients detect equality of
entire original global sections; the coefficient of an original
homogeneous polynomial section is the same original $F/x_i^e$.
\end{lemma}
\begin{proof}
Native restriction of a global module section commutes with the original
frame isomorphisms. The proved original frame transition says on the
whole overlap that $c_i=(x_k/x_i)^e c_k$. Native Proj restriction
naturality identifies these restrictions with the genuine maps between
the original coordinate localizations. Every inverted coordinate is an
actual unit in the original Laurent algebra. Multiplication by $x_i^e$
therefore gives $c_i x_i^e=c_k x_k^e$ in that algebra, using the same
original functions and maps. For the support claim, take the native
homogeneous numerator and denominator of $c_i$. The denominator is a
power of the actual $x_i$ and the numerator has the same original
homogeneous degree as that power. After multiplication by $x_i^e$,
total degree is exactly $e$ and all uninverted exponents remain
nonnegative. Native localization injectivity and clearing establish
these identities with arbitrary coefficient torsion retained.
\uses{lem:fc-original-homogeneous-form-coefficients}
\end{proof}

\begin{lemma}[The entire compatible expression has no negative coordinate]
\label{lem:c0n-ambient-compatible-nonnegative}
\leanok
\lean{EnvelopeProperty.Char0.RelativeProjectiveSectionPolynomial.compatible_singleton_laurent_mem_empty,EnvelopeProperty.Char0.RelativeProjectiveSectionPolynomial.sectionLaurent_mem_empty}
The common expression of the preceding lemma is supported on the whole
original nonnegative degree-$e$ exponent set, including $N=0$.
\end{lemma}
\begin{proof}
For $N\geq1$, fix any coordinate $i$ and choose a different original
coordinate $k$. The common expression equals its expression on chart
$k$, where coordinate $i$ is not inverted. Hence every supported
exponent at $i$ is nonnegative. This holds for every original $i$.
If $N=0$, there is exactly one coordinate and its exponent equals the
entire total degree $e\geq0$, so it is nonnegative as well. Thus the
same conclusion includes projective zero-space over the original ring,
without an implicit nontrivial-coordinate-type assumption.
\uses{lem:c0n-ambient-native-section-laurent}
\end{proof}

\begin{theorem}[Every whole original ambient section has an actual homogeneous lift]
\label{thm:c0n-ambient-section-homogeneous-lift}
\leanok
\lean{EnvelopeProperty.Char0.RelativeProjectiveSectionPolynomial.chartToLaurent_injective,EnvelopeProperty.Char0.RelativeProjectiveSectionPolynomial.exists_homogeneous_polynomial}
For every original $s\in\Gamma(P,\mathcal O(e))$ there exist an actual
$F\in R[x_0,\ldots,x_N]$, homogeneous of degree $e$, such that
\[
 s=\operatorname{homogeneousSection}(F)
\]
in the original coordinate-hyperplane twist. The conclusion uses the
native global section of the same original form, with the original
coefficient ring and the same original frames; no desired lifting map
or comparison certificate is a premise.
\end{theorem}
\begin{proof}
Form the compatible Laurent expression of the original section and
apply the two preceding lemmas. The empty-coordinate reconstruction
then produces the actual original degree-$e$ polynomial $F$.
On each original chart its original polynomial section has coefficient
$F/x_i^e$, by the proved native homogeneous-section coefficient formula.
The constructed equality of Laurent images and invertibility of the
same $x_i^e$ show that this equals $c_i$; actual localization
injectivity brings the equality back to the entire original chart
section ring. The original frame isomorphism is injective, so both
module sections restrict equally on every chart. The original
coordinate opens cover $P$, and its genuine sheaf separatedness proves
equality of the entire original global sections. All arguments in the
previous lemmas include $N=0$, and all maps retain $R$, in particular
when $R=k[t]$.
\uses{lem:c0n-ambient-empty-laurent-polynomial,lem:c0n-ambient-compatible-nonnegative,thm:fc-original-homogeneous-form-section,lem:fc-original-homogeneous-form-frame-coefficient,lem:c0n-polynomial-original-index-reconstruction,lem:c0n-polynomial-localization-laurent-faithful}
\end{proof}

The separate fixed-subscheme restriction-surjectivity theorem is needed
before this ambient lifting result lifts sections on an original closed
projective subscheme. Neither result alone supplies the entire final
characteristic-zero envelope theorem.
